%% annex/p06/axioms-hierarchy.tex --- proofs/sketches for the P6
%% foundations-and-ordinals slice.  One annex item per body result whose
%% status points here; cited results have no annex item.

\begin{annexitem}{annex:p6.axioms.1}{thm:p06-omega-least}{Existence and leastness of \(\omega\)}
\begin{proof}
Every inductive \(J\) contains each selected \(x\), hence \(\omega\subseteq J\).
Since \(I_0\) and every \(J\) contain \(\varnothing\) and are
successor-closed, the defining condition makes \(\omega\) inductive;
Extensionality gives uniqueness.
\end{proof}
\end{annexitem}

\begin{annexitem}{annex:p6.axioms.2}{thm:p06-foundation-no-self-member}{Foundation excludes self-membership}
\begin{proof}
If \(a\in a\), apply Pairing and Foundation to \(\{a\}\).  Its only candidate
minimal member is \(a\), but \(a\in a\cap\{a\}\), a contradiction.
\end{proof}
\end{annexitem}

\begin{annexitem}{annex:p6.axioms.3}{thm:p06-set-recursion}{Set recursion}
\sketchnote{omits the formula-level coding of functions, the construction and
set-existence proof for each transitive closure, and the individual Separation
and Replacement instances; these are supplied by \cite[Chapter~2]{Jech2003}.}
For a set \(x\), form its transitive closure from
\(u_0=x\), \(u_{n+1}=\bigcup u_n\), and
\(\bigcup_{n<\omega}u_n\).  Foundation gives induction on membership.
Define compatible approximations on this transitive closure; Replacement
collects their values, and Union and Separation assemble the resulting graph.
The unique-value hypothesis for \(G\) makes the local approximations coherent.
If two class solutions differed at \(x\), Separation inside the transitive
closure of \(\{x\}\) would give a membership-minimal input where their values
differed.  Their restrictions below it, and therefore their \(G\)-values,
would be equal, a contradiction.  The same set-sized construction on any
domain \(a\) yields the asserted set function
\(F\mathbin{\upharpoonright}a\).
\end{annexitem}

\begin{annexitem}{annex:p6.ordinals.1}{thm:p06-ordinal-trichotomy}{Trichotomy and members of ordinals}
\sketchnote{omits the two elementwise membership calculations in the
proper-subset lemma and cites them to \cite[Chapter~2]{Jech2003}; the repaired
\(\alpha\cap\beta\) route required by R-034 is retained.}
The intersection \(\alpha\cap\beta\) is an ordinal: it is transitive, and
membership strictly well-orders it as a subset of \(\alpha\).  If ordinals
\(\gamma\subsetneq\delta\), let \(\xi\) be the membership-least member of
\(\delta\setminus\gamma\).  Minimality gives \(\xi\subseteq\gamma\), while
linearity on \(\delta\) and transitivity of \(\gamma\) give
\(\gamma\subseteq\xi\); hence \(\gamma=\xi\in\delta\).

If \(\alpha\cap\beta\) were proper in both ordinals, that lemma would put
\(\alpha\cap\beta\) in both, hence in itself, contradicting irreflexivity.
Thus the intersection equals one of \(\alpha,\beta\); equality or the
proper-subset lemma yields the three alternatives, and ordinal transitivity
shows that no two can hold.  Finally, if \(x\in\alpha\), then
\(x\subseteq\alpha\); membership well-orders \(x\), and transitivity of the
membership order on \(\alpha\) makes \(x\) transitive.  Hence \(x\) is an
ordinal.
\end{annexitem}

\begin{annexitem}{annex:p6.ordinals.2}{thm:p06-transfinite-induction}{Transfinite induction}
\begin{proof}
If a counterexample exists, Separation on \(\delta+1\) forms the counterexample
set with least member \(\alpha\).  Every \(\beta<\alpha\) satisfies
\(\varphi\), so progressiveness gives \(\varphi(\alpha)\), contradiction.
\end{proof}
\end{annexitem}

\begin{annexitem}{annex:p6.ordinals.3}{thm:p06-transfinite-recursion}{Transfinite recursion}
\sketchnote{omits the formula-level graph coding and the precise Separation
and Replacement instances at successor and limit stages; these set-existence
details are supplied by \cite[Chapter~2]{Jech2003}.}
For each ordinal \(\gamma\), transfinite induction constructs the unique
function \(f_\gamma\) with domain \(\gamma\) and
\(f_\gamma(\alpha)=G(f_\gamma\mathbin{\upharpoonright}\alpha)\).  Append the
unique value at a successor stage and take the union of the coherent earlier
functions at a limit stage; Replacement collects the earlier graphs.
Transfinite induction also gives
\(f_\beta=f_\gamma\mathbin{\upharpoonright}\beta\) whenever
\(\beta<\gamma\).  Their union therefore defines the asserted class function,
and the same induction proves uniqueness.  Each \(f_\gamma\) is already the
required set restriction.
\end{annexitem}

\begin{annexitem}{annex:p6.hierarchy.1}{thm:p06-hierarchy}{Hierarchy theorem}
\sketchnote{omits the detailed construction of the transitive closure used in
the Foundation argument, the Replacement function assigning stage indices,
and the corresponding Separation instances; these routine coding obligations
are supplied by \cite[Chapters~2--3]{Jech2003}.}
Transfinite induction proves that every \(V_\alpha\) is transitive and that
\(\alpha\leq\beta\) implies \(V_\alpha\subseteq V_\beta\).  If some set lies
in no stage, apply Foundation to the members of the transitive closure of its
singleton that lie in no stage and choose a membership-minimal counterexample
\(x\).  Every member \(y\in x\) has a least stage index \(\alpha_y\);
Replacement collects the ordinals \(\alpha_y+1\).  Their union \(\gamma\) is
a common bound, so all members of \(x\) lie in \(V_\gamma\), and
\(x\in V_{\gamma+1}\), a contradiction.  Finally, ordinal well-ordering gives
the least \(\alpha\) for which \(x\subseteq V_\alpha\).
\end{annexitem}

\begin{annexitem}{annex:p6.hierarchy.2}{thm:p06-mostowski-collapse}{Mostowski collapse}
\sketchnote{omits the construction of the rank function for an arbitrary
well-founded set relation and the exact Separation and Replacement instances
that make the recursive predecessor images sets; see
\cite[Chapter~6]{Jech2003}.}
Well-founded recursion defines
\(\pi(a)=\{\pi(b):bRa\}\).  Replacement makes each predecessor image a set.
An induction on \(R\) uses extensionality of \(R\) to show that \(\pi\) is
injective.  Its range \(B\) is transitive by the defining equation, and that
equation both preserves and reflects the relations, so
\(\pi:(A,R)\to(B,\in)\) is an isomorphism.  A second well-founded induction
shows that any other collapse agrees pointwise with \(\pi\), yielding
uniqueness of the map and its range.
\end{annexitem}
